\documentclass{article}

\begin{document}

\textbf{CHAPTER 3
BASIC CONCEPTS IN 2D}

\hfill

The basic concepts named in the Introduction, and described in
the preceding chapter in a ID context, apply with little modification
in the 2D context which is our main concern in this book. We no
longer have the convenience of a visible graph of the map, however,
because the graph of a 2D map is a 2D surface in a 4D space.
Therefore, we must be satisfied with a frontal view of the 2D
domain of the map, in which we try to visualize as much as
possible.

\hfill

\textbf{3.1 MAPS}

\hfill


As before, by map we mean a continuous function from the
domain to itself, $d:D \to D$. . From now on, the domain will be a twodimensional
subset, usually an open subspace, of the plane. For
example, D might be an open rectangle (that is, not containing its
boundary) or the whole plane. The images and pre images of a
point, the one-to-one property, and noninvertibility are defined as in
2.1. We now consider noninvertible maps, in 2D.

\hfill

Note: The complex number maps familiar from the fractal theories
of Fatou, Julia, Mandelbrot, and others may be regarded as real
2D maps. Thus, they fit in the context of this book.

\hfill

\textbf{3.2 MULTIPLICITIES AND CRITICAL CURVES}

\hfill

In the context of a given map, we define the layer set, zn' as the
set of points having exactly n preimages of rank 1, where n is a natural
number: 0, 1, 2, ... , and so on. Those layer sets that are
nonempty open sets are the multiplicity zones. Points on the boundaries
of the zones, generally, are critical points of the map. In

general, these sets will not exhaust the domain, as there may be
points that have infinitely many preimages of rank 1, and thus do
not belong to zn for any n. Even polynomial maps may have this
problem, but generic (that is, almost all) polynomial maps have the
following nice properties:

\begin{itemize}
\item there are only a finite number of layer sets;
they exhaust the domain;
\item the zones of multiplicity fill almost all of the domain;
\item all layer sets that are not multiplicity zones consist of critical
points, arranged in a set of piecewise smooth curves.
\end{itemize}

A map having these nice properties is called a finitely folded
map, and a curve consisting of critical points is called a critical
curve of rank 1, and is denoted by L.1 The image of a critical curve
of rank 1 is a critical curve of rank 2, denoted Ll , and so on.

Note: For a given map T, $L = T[L_{-1}]$, where $L_{-1}$ is the critical
curve of rank 0, and may be thought of as the set of "coincident preimages"
of points of L.

\hfill

\textbf{3.3 AN EXAMPLE}

\hfill

For example, let D be the entire plane. The polynomial map
$f:D \to D$ defined by f(x, y) = (u, v), where,

\[
u = ax + y
\]
\[
v = b + x^2
\]

is finitely folded. We will set a = - 0.7, and b = 1.0.

Figure 3-1 shows part of the domain of the map, the (x, y) space
D, containing the critical curve of rank 0, $L_{-1}$ which coincides with
the y axis. It divides the domain into two regions, denoted $R_l$ and
$R_2$. Figure 3-2 shows part of the image of the map f, in the (u, v)
space, with the zones of multiplicities zero and two, Zo and Z2' separated
by the critical curve, L. Figure 3-3 shows the two spaces,
superimposed.

There is a folding of the (x, y) space, on the critical curve of rank 0, $L_{-1}$, followed by a nonlinear deformation, a rotation, and a movement to the right, into the space of (u, v). The entire (x, y) space ends up on the zone $Z_2$, with the critical curve $L_{-1}$ moving onto the critical curve, L. 
We may visualize the two regions of the (x, y) space, R 1
and $R_2$ folded onto one another, then distorted and pressed down
onto Z2. Actually, the two regions are mapped onto one.
The motion, visualized in this way, may be reversed. This provides
a method to visualize the action of the inverse mapping as
well. A small area in the (u, v) space on the right, if contained
entirely within Z2' will unfold into two small regions in the (x, y)
space on the left, one in the region R 1 , the other in $R_2$.

As we wish to iterate the map, and to visualize the trajectories,
attractors, and basins, of our discrete dynamical system, it will be
useful (although initially confusing) to superimpose the (u, v) space
on top of the (x, y) space. Then, as a weak substitute for the graphical
method of Koenigs-Lemaray in the ID case, we apply the
motion from Figure 3-1 to Figure 3-2 again and again. The domain,
D, is mapped into itself repeatedly. The curve $L_{-1}$ moves onto the
curve L, which in tum moves onto the curve $L_1$, and so on. This
superimposition is shown in Figure 3-3. This portrait is the basis of
the method of critical curves, which is the main method of this
book.

\hfill

\textbf{3.4 TRAJECTORIES AND ORBITS}

\hfill

In the 2D context, trajectories and orbits are defined exactly as
in 2.3 but here they must be plotted in the two-dimensional domain.
This is particularly appropriate for computer-generated plots. And
in this computational method, the attractors and their basins may
be discovered by experiment. We usually find the critical curve of
rank 0, $L_{-2}$ manually by the standard method of vector calculus
(involving the vanishing of the Jacobian determinant, see Appendix
3), then enter its symbolic description into the computer program,
which can then plot the higher-order iterates. The method of critical
curves is based on experiments such as this.

As in the ID case, there are special kinds of orbits which are
important qualitative features of the dynamics of an iterated map.
First among these are the fixed points, which are unmoved by the
map. The different types of fixed points are defined by the motions
of nearby points. The classification is based on the differential calculus
and linear algebra of two-dimensional spaces, but we will
give here only the results. Excepting certain unusual cases, there are
five kinds of so-called generic fixed points in this classification. The
five types are illustrated in Figures 3-4 to 3-8.

Another special type of orbit of great importance in the qualitative
theory is the periodic orbit, or cyclic orbit, or cycle, which
consists of a finite set of points. The map permutes the points of the
orbit cyclically: If there are n points in the orbit, each of the points
returns to its original position after exactly n iterations of the map.
The number n is called the period of the orbit, which is also called
an n-cycle. A point of an n-cycle is said to have prime period n, and
is also a fixed point of the map iterated n times. Periodic points are
classified according to their type as a fixed point of the iterated
map.

\hfill

\textbf{3.S ATTRACTORS}

\hfill

As in the ID case, there are three types of attractors:

\begin{itemize}
\item static attractors, also called attractive fixed points;
\item periodic attractors, also called cyclic attractors; and
\item chaotic attractors .
\end{itemize}

The static attractors are fixed points which are attractive, that is, the
trajectories of all nearby points are attracted to them. Of the five types of
fixed points illustrated in Figures 3-4 to 3-8, two are attractive, and three
repelling. Periodic attractors are periodic orbits (orbits of trajectories
that cycle around a finite point set) which are attractive.

Chaotic attractors are more complicated sets which are attractive. For
the mathematically inclined, technical definitions are given in Appendix 2.
For others, these concepts will gain meaning through examples later in the
book.

\hfill

\textbf{3.6 BIFURCATIONS}

\hfill

The informal definition of bifurcation, in 20, is the same as in
10. Again, there are subtle and catastrophic bifurcations, and other
distinctions such as local versus global bifurcations. These are best
understood in the examples dissected in detail in the following
chapters.
In the 20 context, a one-parameter family of maps may be displayed
in a response diagram, in which the bifurcations may be
seen and analyzed. This is a 30 plot in which the domain of the
maps, a 20 set, is arrayed vertically, and moved along a horizontal
axis representing the control parameter. In each of these vertical
planes, the portrait of attractors, basins, and boundaries must be
visualized. In practice, this is a challenging task of computer graphics,
and we usually seek a simpler display. The technique we adopt
for this book, which is well accommodated by computer graphic
animation technology and CO-ROM media, is the animated movie.
Thus, we translate the control parameter into the time dimension
and view the domain of the map head on, watching the attractorbasin
portrait adjust itself to a time-changing control parameter.
The method of characteristic curves becomes a strategy for the
analysis of these bifurcation movies. So, on to the exemplary bifurcation
sequences.
BASIC




\end{document}


